% !TEX root = ../hw02.tex
Let
\[
  I_n=\left(0,\frac1n\right),\qquad n\geq1.
\]
Each $I_n$ is a nonempty bounded interval because it contains $1/(2n)$
and is contained in $(0,1)$. Since $1/(n+1)<1/n$,
\[
  I_{n+1}=\left(0,\frac1{n+1}\right)
    \subset I_n.
\]
Thus the intervals are nested, as shown in
Figure~\ref{fig:hw02:open-intervals}.
\begin{marginfigure}
  \centering
  % !TEX root = ../../problems/hw02.tex
% (0,1/n), n = 1,2,3,..., drawn on the same horizontal scale.
\begin{tikzpicture}[interval diagram]
  \foreach \n/\endpoint in {1/1,2/{\tfrac12},3/{\tfrac13}} {
    \coordinate (left) at (0,{1-\n});
    \coordinate (right) at ({3.6/\n},{1-\n});
    \draw (left) -- (right);
    \node[interval open] at (left) {};
    \node[interval open] at (right) {};
    \node[interval label,above=3pt] at (right) {$\endpoint$};
  }
  \node[interval label,above=3pt] at (0,0) {$0$};
  \node at (1.8,-2.9) {$\vdots$};
\end{tikzpicture}

  \caption{The nested intervals $I_n=(0,1/n)$ exclude $0$.}
  \label{fig:hw02:open-intervals}
\end{marginfigure}

To prove that their intersection is empty, let $x\in\R$.
If $x\leq0$, then $x\notin I_1$. If $x>0$, the Archimedean
property gives a positive integer $n>1/x$.\sidenote{\textit{Archimedean property.}
  For every $M\in\R$, there is a positive integer $n>M$.}
Then $1/n<x$ and $x\notin I_n$.
Hence every real number is excluded from at least one interval,
and therefore
\[
  \bigcap_{n=1}^{\infty}I_n=\varnothing.
\]

The nested interval property guarantees a nonempty intersection
for nested nonempty bounded closed intervals in $\R$. Here the
missing hypothesis is \emph{closedness}. Every finite intersection
is nonempty, but the only point common to all the closures
$[0,1/n]$ is $0$, which is excluded from every $I_n$.
